% !TEX program = xelatex
% Шаблон научной статьи НИТУ «МИСиС»
% Заполните блоки ниже и скомпилируйте: xelatex article.tex

\documentclass{misisarticle}

\usepackage{fontspec}
\setmainfont{Times New Roman}[
  SmallCapsFont = {Times New Roman},
  SmallCapsFeatures = {Letters=SmallCaps, FakeBold=1.2}
]

% --- Классификационные коды (оставьте пустым {}, если не требуются) ---
\grnti{28.23.37, 84.37.03, 47.23.15}
\asjc{2208, 3313}
\oecd{2.2, 6.2}

% --- Название ---
\title{Использование UWB-антенн для задач точного позиционирования трамвайных вагонов}

% --- Авторы: \misisauthor{ФИО}{аффилиация}{город, страна}{email} ---
\misisauthor{М.~А.~Шульгин}{кафедра инженерной кибернетики НИТУ «МИСиС»}{Москва, Россия}{m2311717@edu.misis.ru}
\misisauthor{Р.~Н.~Садеков}{кафедра инженерной кибернетики НИТУ «МИСиС»}{Москва, Россия}{r.sadekov@misis.ru}

% --- Аннотация и ключевые слова ---
\abstract{%
  В статье рассматривается проблема обеспечения надежного и высокоточного позиционирования трамвайных вагонов в условиях плотной городской застройки и на подземных участках, где традиционные ГНСС-системы неэффективны. Предложен подход к применению сверхширокополосной (UWB) радиолокации на основе методов ToF и TDoA. Особое внимание уделено оптимизации конструкции и пространственного размещения UWB-антенн на корпусе вагона для минимизации влияния металлических конструкций и контактной сети. Описана методика валидации решения путем сопоставления с эталонными данными референсной системы.

}

\keywords{%
  сверхширокополосная радиолокация (UWB), точное позиционирование, трамвайный транспорт, ГНСС-отказоустойчивые среды, метод Time of Flight (ToF), метод Time Difference of Arrival (TDoA).
}

\begin{document}

\maketitle

% ============================================================
\section{Введение}
Современные системы автоматической диспетчеризации и обеспечения безопасности движения наземного электротранспорта требуют непрерывного и высокоточного определения координат подвижного состава. Традиционные глобальные навигационные спутниковые системы (ГНСС), такие как GPS и ГЛОНАСС, демонстрируют низкую эффективность в условиях плотной городской застройки («городские каньоны») и на подземных участках трамвайных маршрутов. В этих средах сигналы ГНСС подвержены сильному экранированию и многолучевому распространению (multipath), что приводит к значительным погрешностям в определении местоположения или полной потере навигационного поля \cite{huang2024design}. В связи с этим актуальной научно-технической задачей становится разработка и внедрение альтернативных или комплементарных технологий позиционирования, способных обеспечить субметровую точность в ГНСС-отказоустойчивых средах.


% --- 2. Анализ существующих решений ---
\section{Анализ существующих решений}
Проблема надежной локализации транспорта в условиях отсутствия или деградации ГНСС-сигнала является одной из ключевых в развитии интеллектуальных транспортных систем. Исторически потенциал сверхширокополосных (UWB) технологий для решения этой задачи был продемонстрирован в работах по гибридному позиционированию. В частности, в исследовании \cite{fernandez2007application} обоснована эффективность совместного использования UWB и GPS для бесшовной локализации транспортных средств при переходе из открытых пространств в закрытые. 

В современных работах акцент смещается на алгоритмическое совершенствование UWB-систем. Так, авторы \cite{huang2024design} предлагают новые алгоритмы фильтрации, позволяющие существенно повысить точность позиционирования именно в средах, недоступных для GPS, за счет подавления многолучевых эффектов. Методологическую базу для оценки таких систем в сложных условиях формируют исследования, подобные \cite{queralta2020uwb}, где, несмотря на фокус на беспилотных летательных аппаратах, детально характеризуются особенности распространения UWB-сигналов и формируются репрезентативные датасеты для ГНСС-отказоустойчивых сред, что применимо и к наземному транспорту в условиях «городских каньонов».

В контексте именно трамвайного движения в отечественной практике наибольшее развитие получили визуально-геоинформационные методы и методы сопоставления с цифровыми картами. В работах \cite{ali2025, guzhva2023} исследуется оценка точности местоположения трамвая в условиях высотной застройки с использованием сегментированных изображений и данных карт. Альтернативный подход, основанный на использовании монокамеры и дорожных знаков с геодезической привязкой, предложен в \cite{bikmaev2020}. 

Несмотря на высокую потенциальную точность, оптические и визуально-геоинформационные методы имеют фундаментальные ограничения для систем безопасности критического уровня. Они крайне зависимы от внешних факторов: освещенности, погодных условий (снег, туман, загрязнение объектива), а также требуют значительных вычислительных ресурсов для обработки видеопотока в реальном времени. Кроме того, эти методы часто оказываются неработоспособными на подземных участках трамвайных маршрутов, где визуальные ориентиры отсутствуют или ненадежны.

Таким образом, выявляется исследовательский пробел: отсутствие комплексных решений по внедрению UWB-локации именно для трамвайного подвижного состава, учитывающих специфику его конструкции. Большинство существующих разработок не учитывают сильное экранирование сигнала массивным металлическим кузовом вагона и специфические электромагнитные помехи от тяговой контактной сети. Это обуславливает необходимость разработки специализированного подхода к оптимизации размещения и настройки UWB-антенн на трамвае для обеспечения субметровой точности и высокой отказоустойчивости системы.

% --- 3. Постановка задачи ---
\section{Постановка задачи}
Целью данного исследования является обеспечение надежного и высокоточного позиционирования трамвайных вагонов в условиях плотной городской застройки и подземных участков, где традиционные системы ГНСС демонстрируют низкую эффективность из-за многолучевости и экранирования сигнала. 

Достижение цели направлено на улучшение показателей точности, отказоустойчивости и быстродействия систем автоматической диспетчеризации и безопасности движения подвижного состава. Для решения поставленной задачи предлагается способ применения сверхширокополосной (UWB) радиолокации на основе времяпролетного метода (Time of Flight, ToF) и дифференциального метода (Time Difference of Arrival, TDoA). Ключевым аспектом исследования является оптимизация конструкции и пространственного размещения UWB-антенн на корпусе вагона для минимизации влияния массивных металлических конструкций кузова и электромагнитных помех от контактной сети. Эффективность разработанного решения предполагается проверять путем проведения натурных испытаний на действующем трамвайном маршруте (или физическом макете) с сопоставлением полученных координат с эталонными данными высокоточной референсной системы.

% --- 4. Описание методов ---
\section{Описание предлагаемого метода}
Решение поставленной задачи базируется на комплексном подходе, включающем выбор физической основы позиционирования, инженерную оптимизацию аппаратной части и строгую методологию валидации.

\subsection{Принцип работы UWB-локализации (ToF и TDoA)}
Сверхширокополосные сигналы обладают высокой временн\'{о}й разрешающей способностью, что позволяет им эффективно разделять прямой сигнал и многолучевые отражения. В предлагаемой системе применяются два взаимодополняющих метода:
\begin{itemize}
    \item \textbf{Метод ToF (Time of Flight):} измеряет абсолютное время прохождения радиосигнала от передатчика на вагоне до стационарных якорных узлов (anchors), установленных вдоль пути. Требует высокой синхронизации часов, но обеспечивает высокую точность на открытых участках.
    \item \textbf{Метод TDoA (Time Difference of Arrival):} измеряет разницу во времени прибытия сигнала от передатчика вагона к нескольким синхронизированным между собой якорным узлам. Этот метод менее требователен к синхронизации часов на самом вагоне и лучше подходит для плотной городской среды, так как позволяет вычислять положение по гиперболам.
\end{itemize}

\subsection{Оптимизация размещения антенн на корпусе вагона}
Главной технической проблемой при внедрении UWB на трамвае является сильное экранирование и переотражение сигналов от металлического кузова, а также интерференция с проводами контактной сети. Для минимизации этих эффектов предлагается следующий методологический подход:
\begin{enumerate}
    \item \textbf{Антенное разнообразие (Diversity):} Использование нескольких UWB-антенн, разнесенных по разным углам или сторонам вагона (например, на крыше и на уровне рамы тележки). Это позволяет системе динамически выбирать антенну с наилучшим отношением сигнал/шум в зависимости от ориентации вагона относительно базовых станций.
    \item \textbf{Экранирование и направленность:} Применение антенн с направленной диаграммой направленности, ориентированной в сторону базовых станций (вдоль пути), и использование поглощающих экранов для блокировки сигналов, идущих от контактной сети сверху.
    \item \textbf{Моделирование:} Предварительное электромагнитное моделирование (например, в CST Studio Suite или HFSS) геометрии вагона для определения «мертвых зон» и оптимальных точек крепления антенных модулей.
\end{enumerate}

\subsection{Методика валидации решения}
Для подтверждения эффективности предложенного метода планируется проведение натурных испытаний. Методика включает:
\begin{itemize}
    \item Развертывание сети UWB-якорей на выбранном участке трамвайного маршрута (включая проблемные зоны: тоннели, участки с плотной застройкой).
    \item Оснащение трамвайного вагона прототипом UWB-системы с оптимизированным размещением антенн.
    \item Использование высокоточной референсной системы (например, тахеометрической съемки на закрытых участках или RTK-GNSS с инерциальной навигационной системой (INS) на открытых участках) для получения эталонной траектории движения (Ground Truth).
    \item Статистическое сравнение координат, полученных от UWB-системы, с эталонными данными с расчетом среднеквадратичной ошибки (RMSE), максимальной погрешности и процента времени доступности сигнала.
\end{itemize}

% --- 5. Список литературы ---
\begin{thebibliography}{9}

\bibitem{fernandez2007application}
J. A. Fernandez-Madrigal, E. Cruz-Martin, J. Gonzalez, C. Galindo and J. L. Blanco, "Application of UWB and GPS technologies for vehicle localization in combined indoor-outdoor environments," \textit{2007 9th International Symposium on Signal Processing and Its Applications}, Sharjah, United Arab Emirates, 2007, pp. 1-4, doi: 10.1109/ISSPA.2007.4555416.

\bibitem{huang2024design}
Y. Huang, B. Cao, and A. Wang, "Design a novel algorithm for enhancing UWB positioning accuracy in GPS-denied environments," \textit{Sci. Rep.}, vol. 14, Art. no. 23895, 2024, doi: 10.1038/s41598-024-74773-y.

\bibitem{queralta2020uwb}
J. P. Queralta, C. Martínez Almansa, F. Schiano, D. Floreano and T. Westerlund, "UWB-based System for UAV Localization in GNSS-Denied Environments: Characterization and Dataset," \textit{2020 IEEE/RSJ International Conference on Intelligent Robots and Systems (IROS)}, Las Vegas, NV, USA, 2020, pp. 4521-4528, doi: 10.1109/IROS45743.2020.9341042.

\bibitem{ali2025}
Али Б., Садеков Р.Н. Оценка точности позиционирования трамвая в кривых на основе данных карты и сегментированных изображений // Научно-технический вестник информационных технологий, механики и оптики. 2025. Т. 25, № 4. С. 771-779.

\bibitem{guzhva2023}
Оценка точности местоположения трамвая в условиях высотной застройки по данным визуально-геоинформационных систем / Н. С. Гужва, Б. Али, К. С. Бакулев [и др.] // XXX Юбилейная Санкт-Петербургская международная конференция по интегрированным навигационным системам : сборник материалов конференции, Санкт-Петербург, 29–31 мая 2023 года. – Санкт-Петербург: Концерн "ЦНИИ "Электроприбор", 2023. – С. 16-21. – EDN: VCHMDX.

\bibitem{bikmaev2020}
Бикмаев Р. Р., Полукаров А. А., Садеков Р. Н. Определение местоположения наземного транспортного средства с использованием монокамеры и дорожных знаков с геодезической привязкой // XXVII Санкт-Петербургская международная конференция по интегрированным навигационным системам : Сборник материалов, Санкт-Петербург, 25 мая – 05 июня 2020 года. – Санкт-Петербург: Концерн "ЦНИИ "Электроприбор", 2020. – С. 341-345. – EDN: TDXGEK.

\end{thebibliography}

\end{document}